\documentclass[11pt,twoside,openright]{book}

\usepackage[T1]{fontenc}
\usepackage[utf8]{inputenc}
\usepackage{amsmath,amssymb,amsthm,mathtools,bm}
\usepackage{mathpazo}
\usepackage{microtype}
\usepackage[inner=1.35in,outer=1.1in,top=1.25in,bottom=1.35in]{geometry}
\usepackage{booktabs}
\usepackage{fancyhdr}
\usepackage[hidelinks]{hyperref}
\usepackage{cleveref}

\setlength{\parindent}{1.4em}
\setlength{\parskip}{0pt}
\linespread{1.07}
\emergencystretch=3em

\pagestyle{fancy}
\fancyhf{}
\fancyhead[LE]{\small\thepage\quad\textsc{A Cinema of Individual Editions}}
\fancyhead[RO]{\small\nouppercase{\leftmark}\quad\thepage}
\renewcommand{\headrulewidth}{0.3pt}
\fancypagestyle{plain}{\fancyhf{}\fancyfoot[C]{\small\thepage}\renewcommand{\headrulewidth}{0pt}}

\theoremstyle{plain}
\newtheorem{theorem}{Theorem}[chapter]
\newtheorem{proposition}[theorem]{Proposition}
\newtheorem{lemma}[theorem]{Lemma}
\newtheorem{corollary}[theorem]{Corollary}
\theoremstyle{definition}
\newtheorem{definition}[theorem]{Definition}
\newtheorem{example}[theorem]{Example}
\theoremstyle{remark}
\newtheorem{remark}[theorem]{Remark}

\newcommand{\Capac}{\mathrm{Cap}}
\newcommand{\Rob}{\mathrm{Rob}}
\newcommand{\Inv}{\mathrm{Inv}}
\DeclareMathOperator{\loss}{loss}

\title{\Huge\bfseries A Cinema of Individual Editions\\[1.2ex]
\large Chapter Eighteen: Open Problems\\[0.4ex]
\large Working Draft}
\author{Flyxion\\[0.3ex]\small Independent Researcher}
\date{November 2026}

\begin{document}

\frontmatter
\maketitle
\tableofcontents
\mainmatter
\setcounter{chapter}{17}

\chapter{Open Problems}
\label{ch:open}

This chapter collects what the book leaves undecided. It is ordered by what each problem would unlock, and it keeps two kinds of question apart. The mathematical problems are about the models of the preceding chapters and would be settled by proof or by computation. The empirical problems are about viewers, footage, and productions, and no amount of mathematics settles them, however many tables the earlier chapters contain. The chapter introduces no new counts.

It also has to say what the book does not show. It does not show that a cinema of individual editions is desirable, feasible, or something that anyone should build. Every capacity in it is a property of a toy scene of five slots, every tolerance is a stand-in for a threshold that nobody has measured, and the rate of Chapter~1 has a denominator that has never been computed. The arithmetic is exact, and it is exact about scenes that were chosen so that it could be. The book does not support its own premise, and the problems below are the reasons.

\section{Mathematical Problems}\label{sec:math18}

\paragraph{Exact covering numbers.} Chapter~2 gives the packing numbers of the corrected pools exactly and the covering numbers as ranges where the solver stopped: for production A at tolerance one the covering number lies between $25$ and $31$ and was not resolved. A covering is the smallest set of editions such that every edition is within the tolerance of one of them, and so it is the number of editions needed to represent the pool at that resolution. It is related to the packing number by the inequalities of Chapter~2 and is not otherwise its dual: a packing governs the simultaneous separation of viewers, and a covering governs representation within a tolerance. It is the quantity that a compression claim would use. Closing the gap is a finite computation on a pool of $238$ tuples, and a structural argument that gave the covering number for a class of pools would be better than the computation.

\paragraph{Rules that reach past neighbors.} Counting and unranking in Chapters~13 and~16 depend on the admissibility of a tuple being decided by adjacent pairs, or by a bounded frontier as in Chapter~5, where the frontier widths of the lawn scene are one, two, and three. A rule on the whole edition has no table and sends a generator back to rejection. The open question is which families of rules keep counting and unranking polynomial, and which make counting hard. A dichotomy of the kind known for constraint satisfaction problems is the form of answer one would want, and none is claimed.

\paragraph{Computing robust capacity.} Chapter~9 defines robust capacity as the minimum over the amendments in a family of the capacity left after each, with a separate allocation for each, and invariant capacity as the capacity of the intersection. Both were computed by enumeration on the lawn scene. The complexity of computing them for a family of amendments given as rules, without enumerating the pool, is open, as is which structured families of amendments admit a useful bound between the two. No bound holds for arbitrary families: two survivor sets of size $k$ with one class in common have robust capacity $k$ and invariant capacity $1$, and with no class in common invariant capacity $0$. In the lawn scene they differ by a factor of about one and a half for production A, $351$ against $238$, and a general bound would say when anticipating a correction is as good as waiting for it.

\paragraph{Repair under repeated amendment.} Chapter~10 treats repair as a minimum-cost assignment of invalid viewers to free valid editions, and the radius, the cost, and the number of moves are three different objectives. In production B the repair of minimum total movement has radius three, while a repair of radius two exists, and in sequence the cost is $28$ against $27$ for the one-shot repair. The open problem is the online one: when amendments arrive without notice and each must be repaired before the next is known, how much more than the offline optimum must be paid, in cost and in the number of times a viewer is moved. No bound is proved, and the lawn scene shows only that the gap can be nonzero.

\paragraph{Worldwide footage with local and shared slots.} Chapter~11 proves that locale-by-locale allocation is sound for tolerances below the number of slots that carry a local trace, and gives lower bounds, from the Singleton bound, on the classes per slot that a locale needs. The bounds for an audience of $238$ were not checked for attainment, and the problem of how to divide footage between local and shared slots so as to minimize total footage, subject to the Hall condition across locales and a separation tolerance, is posed and not solved. Codes over mixed alphabets with some coordinates shared between several codes are the likely tool.

\paragraph{Designed schedules.} The cyclic schedule of Chapter~13 gives every audience distinct editions in every round and every viewer the whole pool without repeat, and it guarantees exact inequivalence only. Chapter~16 shows that a generator that selects at random reaches roughly seven tenths to nine tenths of the optimum packing in the cases tried and seldom reaches the optimum. What is wanted is a schedule that gives simultaneous separation by a tolerance and, for each viewer, a sequence of mutually separated editions for the whole novelty horizon. The difficulty is that the pool is a product with tuples removed by rules, and the algebraic constructions that give optimal codes assume a full product. Constructions inside constrained pools are open.

\paragraph{The class of order-independent updates.} Chapter~14 gives two updates, intersection and revision, proves that the final state is independent of order when the committed sets have a common intersection and dependent on order for two films whose sets are disjoint, and shows on the engine scene that the openness at first viewing can depend on order. Families between the two cases, such as films that overlap pairwise without a common intersection, are not characterized, and neither are the updates that are order-independent in the final state or in the whole path. It would say which models of understanding admit a book's worth of films read in any order, and which do not.

\paragraph{Exclusion beyond half the tolerance.} The bound of Chapter~15 removes at most one packing point per ball when the neighborhood radius is at most half the tolerance, and the covering form of it gives $M_\varepsilon-c$ for a removed region covered by $c$ balls of radius $\varepsilon/2$. The survivor counts show that a larger radius can take the whole pool quickly. What is open is how to obtain a useful value of $c$ from the representation and the exclusion policy without enumerating the pool, as is the complexity of deciding whether a policy retains every class of the grammar above a stated fraction.

\paragraph{When local sampling is fair.} Chapter~16 shows that a left-to-right generator without lookahead is uniform over its output if and only if the product of option counts is the same for all admissible tuples. Which families of rules have this property is not characterized. The rules of the lawn scene do not, and a result that said which rules do would say when a table is unnecessary.

\paragraph{Anonymity against capacity.} Chapter~12 shows that a ledger of unique editions is a ledger of near-identifiers. A production could promise instead that each edition is shared by at least $k$ viewers, which relaxes the capacity requirement from the audience to the audience over $k$ and restores an anonymity set of $k$. The trade between the individuality that the book promises and the anonymity that its ledger needs is a one-parameter family that the book states and does not study. What the least pool is for a given $k$ and tolerance, and how far the promise of individuality survives at $k>1$, is open.

\section{Empirical Problems}\label{sec:emp18}

\paragraph{The storage cost.} The rate of Chapter~1 is $\log_2K/B$, with $B$ the cost of the footage, rules, and metadata needed to reconstruct the editions. The book computes $K$ and never $B$. The claim that individual editions are a better compression of film than a single edition is therefore not tested by anything here, and it would be tested by the measurement of $B$ for one real production, with the perceptual distinctions among its editions established, and by a baseline: the storage needed by two representations to deliver the same catalogue, or the same perceptual coverage at a stated tolerance, measured at the same revision. This is the first problem in the order of what it would unlock, because the rate is the justification of the framing and it has a missing denominator.

\paragraph{The thresholds of perceptual equivalence.} Chapter~2 defines three pseudometrics, for detection, comprehension, and affect, and the tolerances at which two editions count as the same are unmeasured. Every packing number, every novelty horizon, and every count of distinguishable editions in the book depends on them. A measurement would give a threshold for each kind of difference and each kind of viewer, and would say whether the Hamming dissimilarity on class tuples that the book uses is anywhere near the measured one.

\paragraph{Local trace.} Chapter~11 assumes that local footage at a slot differs perceptibly between locales and that shared slots do not. Whether it does is a fact about footage and viewers. A comparison of takes across locales, by viewers who are asked whether they differ, would settle which slots carry local trace and so how far locale-by-locale allocation is sound.

\paragraph{Shared viewing.} The exchange of Chapter~1, in which two viewers both say that the same character missed a strip of lawn, one from a close shot and one from inferring it from the finished lawn, is the reason the book gives for wanting editions to be watched together. Whether viewers who have seen different editions can have such an exchange, how much of the film they can discuss, and what they take to be the film, is not known. It is the question on which the practical value of the whole programme depends.

\paragraph{Openness as experienced.} Chapters~3, 14 and~17 count the readings that an edition leaves open and treat the order of films as changing how long a question stays open. Whether viewers experience openness as the size of a set of readings, whether a revealed reading is recognized as the one that was left open, and whether a viewer who is settled at the first viewing sees a different film, are questions for a reading task of the kind Study One of Chapter~17 describes.

\paragraph{Recognition and the tenth year.} Chapter~14 shows that in its model a later release can only lose openness, and leaves open whether a reading that requires a premise is recognized differently by viewers with and without it. The experiment that would settle it holds the edition fixed, varies the viewer's prior films, and asks which readings are reported. Without it the question whether a release in the tenth year of a viewer's life does something that a release in the first would miss has no answer.

\paragraph{Recovery windows and the dialogue reference.} Chapter~6 and Study Three of Chapter~17 take the occupation record at its threshold of five seconds and its reference of $-34.554$ LUFS, and show that a shift of $0.44$ decibels in either direction leaves the passage's three windows unchanged, although the upward shift reduces the film's nine blocks to seven, and that a downward shift to the centre-channel value of $2.33$ decibels leaves two windows of five seconds with zero slack, which differs from having no windows. Whether a window of five seconds at or below the reference is experienced as recovery, which of three references is the level of dialogue, and whether the saturating response of Chapter~6, which is a stipulated model, describes viewers at all, are listening questions.

\paragraph{How corrections arrive.} The ledgers of Chapters~10, 12 and~17 assume corrections that arrive as discrete events with a known rule, and the claim rates of Chapter~7 are assumed. How real reports arrive, how many are true, how long adjudication takes, and whether the amendments of a real production have the shape of the three of the lawn scene, would come from a production's own log replayed against a ledger of its audience.

\paragraph{What viewers disclose.} The anonymity sets of Chapter~12 are computed at the class level, from an observer who knows the class of a viewer's edition at some slots. How much of an edition viewers reveal in talking about a film, and to whom, is unknown, and so is the number of slots an observer typically has. The count says what an observation of $s$ slots leaves, and not how often one occurs.

\paragraph{Whether it makes viewing together easier.} The larger aim of the programme is that individual editions might one day lessen the need to watch films all together, and make a society more interesting and easier to move around in. The book contains no evidence about any society. Whether individual editions make shared experience richer, or poorer, or neither, depends on how people who have seen different films talk to each other, and that is learned from people and not from capacity.

\section{What Each Would Unlock}\label{sec:order18}

The problems are not equally urgent. The measurement of $B$ and the measurement of the thresholds of equivalence come first, because the rate and every count of distinguishable editions are conditional on them, and until they are made the book's central quantities are the numbers of a toy. The question of shared viewing comes next, since it decides whether the programme has a purpose that survives the toy. The design problems, for schedules inside constrained pools, for footage across locales, and for anticipated allocation, come after, because they are only worth solving at the thresholds that the measurements give. The remaining mathematical problems, on covering numbers, on rules beyond neighbors, on online repair, and on the class of order-independent updates, are independent of the measurements and can be pursued at once. The experiments on openness, recognition, and recovery windows are the ones that connect the model of meaning to viewers, and they unlock the chapters on sequence and sound.

None of them is begun here. The book states what it counts, what the counts assume, and what would have to be found out, and it does not claim to have found any of it out.

\section{What the Chapter Establishes and What It Assumes}\label{sec:limits18}

The chapter establishes only a list, and the list is incomplete: it contains the problems that the author of the book met while writing it, and a reader will meet others. The order is a judgment about unlocking, and a different judgment would give a different order. The separation of the mathematical from the empirical is cleaner on the page than it is in practice, since several of the mathematical problems, such as the choice of footage across locales, are posed with parameters that only the empirical problems can supply.

No script accompanies this chapter, since it contains no count.

\begin{thebibliography}{99}

\bibitem{flyxion2026gf}
Flyxion.
\newblock \emph{Generative Film: Categories, Sheaves, and Types for Cinema}.
\newblock Unpublished monograph, October 2026.

\end{thebibliography}

\end{document}
